\subsubsection{Xấp xỉ Laplace [Nâng cao]}

Xấp xỉ Laplace là một phương pháp xấp xỉ phân phối hậu nghiệm của mô hình
Logistic Regression bằng một phân phối Gaussian, với tâm tại nghiệm MAP
(Maximum A Posteriori). Ý tưởng cốt lõi là khai triển bậc hai của log-posterior
quanh điểm cực đại $\boldsymbol{w}_{MAP}$.

Cụ thể, ta có posterior: \begin{equation}
  p(\boldsymbol{w} \mid \mathcal{D}) \propto p(\mathcal{D} \mid \boldsymbol{w})
  p(\boldsymbol{w}).
\end{equation}

Ta xấp xỉ log-posterior bằng khai triển Taylor bậc hai tại
$\boldsymbol{w}_{MAP}$: \begin{equation}
  \ln p(\boldsymbol{w} \mid \mathcal{D}) \approx \ln p(\boldsymbol{w}_{MAP}
  \mid \mathcal{D}) - \frac{1}{2} (\boldsymbol{w} -
  \boldsymbol{w}_{MAP})^\intercal \boldsymbol{A} (\boldsymbol{w} -
  \boldsymbol{w}_{MAP}),
\end{equation} trong đó ma trận Hessian âm được định nghĩa là: \begin{equation}
  \boldsymbol{A} = - \nabla \nabla \ln p(\boldsymbol{w} \mid \mathcal{D})
  \Big|_{\boldsymbol{w} = \boldsymbol{w}_{MAP}}.
\end{equation}

Từ đó, posterior được xấp xỉ bởi phân phối Gaussian: \begin{equation}
  q(\boldsymbol{w}) = \mathcal{N}(\boldsymbol{w} \mid \boldsymbol{w}_{MAP},
  \boldsymbol{A}^{-1}),
\end{equation} trong đó, $\boldsymbol{A}^{-1}$ đóng vai trò là ma trận hiệp
phương sai, mô tả độ bất định của các tham số mô hình.

\paragraph{Bất định của đường biên quyết định}

Trong Logistic Regression, đường biên quyết định được xác định bởi:
\begin{equation}
  \boldsymbol{w}^\intercal \boldsymbol{x} = 0.
\end{equation} Với xấp xỉ Laplace, do $\boldsymbol{w}$ là một biến ngẫu nhiên
Gaussian, biên quyết định trở thành một phân phối thay vì một đường cố định.

Gọi phương sai của logit tại một điểm $\boldsymbol{x}$ là: \begin{equation}
  \sigma^2(\boldsymbol{x}) = \boldsymbol{x}^\intercal \boldsymbol{A}^{-1}
  \boldsymbol{x}.
\end{equation} Khi đó, ta có thể biểu diễn vùng bất định của biên quyết định
xấp xỉ bởi: \begin{equation}
  \boldsymbol{w}_{MAP}^\intercal \boldsymbol{x} \pm 2 \sigma(\boldsymbol{x}).
\end{equation} Vùng này tương ứng với khoảng tin cậy xấp xỉ $95\%$, thể hiện
mức độ dao động của ranh giới phân lớp do bất định tham số.

\figurename~\ref{fig:laplace-decision-boundary-sigma} cho thấy vùng bất định
của đường biên quyết định dựa trên mô hình hồi quy logistic với tập dữ liệu
được sinh trong phần Probit (xem biểu diễn dữ liệu này tại
\figurename~\ref{fig:probit-data}).

\begin{figure}[htbp]
  \centering
  \includegraphics[width=0.9\textwidth]{laplace-decision-boundary-sigma}
  \caption{Vùng bất định của đường biên quyết định dựa trên xấp xỉ Laplace}
  \label{fig:laplace-decision-boundary-sigma}
\end{figure}

Ta thấy vùng bất định trong hình khá hẹp. Điều này cho thấy mô hình tự tin cao
và biên quyết định ổn định. Có thể do dữ liệu tách biệt tốt.
